\chapter{Drug Interactions}

\section*{Definition and Overview}
A drug interaction occurs when one medicine, food, drink, or supplement changes the effect of another medicine. Interactions can make a medicine stronger or weaker, cause new side effects, or make a medicine ineffective. They can occur between prescription medicines, over-the-counter medicines, herbal and complementary products, and with alcohol and certain foods. Many interactions are harmless or manageable, but some are serious and can cause hospitalisation or death. Being aware of interactions, telling health professionals about all medicines, and checking before adding new ones reduces the risk. This chapter explains the types of interactions and how to avoid them.

\section*{Epidemiology}
Drug interactions are a significant cause of harm, especially in older people, who often take many medicines. Around one in five older Australians takes five or more medicines, and the risk of interactions rises with each additional medicine. A substantial proportion of hospital admissions in older people are related to medicines, including interactions. Common culprits include blood thinners, blood pressure medicines, anti-inflammatory drugs, and medicines for diabetes, along with alcohol and some foods.

\section*{Aetiology and Pathophysiology}
\begin{itemize}
    \item \textbf{Pharmacokinetic interactions:} one medicine changes how another is absorbed, distributed, metabolised, or eliminated
    \item \textbf{Pharmacodynamic interactions:} two medicines have similar or opposing effects, changing the overall result
    \item \textbf{Enzyme effects:} some medicines speed up or slow down the liver enzymes that process other medicines
    \item \textbf{Food interactions:} grapefruit juice, alcohol, and some foods affect medicine metabolism
    \item \textbf{Complementary products:} herbal and supplement products can interact with medicines
    \item \textbf{Risk factors:} multiple medicines, older age, kidney or liver disease, and certain medicine classes
\end{itemize}

\section*{Clinical Features}
\begin{itemize}
    \item Increased side effects, such as bleeding, drowsiness, or low blood pressure
    \item Reduced effectiveness of a medicine
    \item Unexpected symptoms after starting a new medicine
    \item Interactions may be immediate or develop over time
    \item Common examples: warfarin with antibiotics or NSAIDs increases bleeding risk; grapefruit juice with some statins increases side effects
    \item Alcohol interacts with many medicines, increasing sedation and liver effects
\end{itemize}

\section*{Differential Diagnosis}
\begin{itemize}
    \item Interactions versus side effects of a single medicine
    \item Interactions versus the effects of an underlying illness
    \item Medicine-food interactions versus medicine-medicine interactions
    \item The mechanism of the interaction, which guides management
\end{itemize}

\section*{When to Seek Medical Care}
Seek medical advice before starting, stopping, or changing any medicine, and before taking complementary products with prescription medicines. Tell every doctor, pharmacist, and health professional about all medicines you take, including over-the-counter products.

\textbf{Call 000 (or local emergency services) for:}
\begin{itemize}
    \item Unexpected bleeding, severe drowsiness, fainting, or confusion
    \item Difficulty breathing or a severe allergic reaction
    \item Irregular heartbeat or chest pain
    \item Seizures or collapse
\end{itemize}

\section*{Diagnosis and Investigations}
\begin{itemize}
    \item \textbf{Medication review:} a complete list of all medicines, including over-the-counter and complementary products
    \item \textbf{Interaction checking:} pharmacists and doctors use software and knowledge to identify interactions
    \item \textbf{Blood tests:} monitoring of medicines with narrow safety margins, such as warfarin
    \item \textbf{Review of symptoms:} assessment of possible interaction effects
    \item \textbf{Home Medicines Review:} a pharmacist can review medicines at home for people at risk
\end{itemize}

\section*{Management}
\begin{itemize}
    \item \textbf{Adjusting doses:} changing a dose to manage an interaction
    \item \textbf{Timing separation:} taking interacting medicines at different times
    \item \textbf{Switching medicines:} choosing an alternative without the interaction
    \item \textbf{Stopping the culprit:} discontinuing a medicine or product causing the interaction
    \item \textbf{Monitoring:} closer monitoring of blood tests and symptoms
    \item \textbf{Avoiding alcohol and grapefruit:} where interactions require it
    \item \textbf{Professional advice:} checking with a pharmacist before adding any new product
\end{itemize}

\section*{Complications}
\begin{itemize}
    \item Serious bleeding, including stomach bleeding from NSAIDs with blood thinners
    \item Dangerous sedation and breathing problems from combined sedatives
    \item Kidney or liver damage
    \item Uncontrolled blood pressure or blood sugar
    \item Toxicity of medicines such as lithium or digoxin
    \item Hospitalisation and death in severe cases
\end{itemize}

\section*{Prognosis and Outcomes}
Most drug interactions can be prevented or managed with awareness and professional review. People who keep an up-to-date list of all medicines and consult pharmacists and doctors before changes have far lower risks. Interaction checking is routine in pharmacies, and most interactions are identified and managed safely. With good communication, the risk of serious harm is small.

\section*{Prevention and Self-Care}
\begin{itemize}
    \item Keep a complete, current list of all medicines and supplements
    \item Show the list to every doctor and pharmacist
    \item Use one pharmacy for all prescriptions
    \item Read the labels and information provided with medicines
    \item Ask before adding any over-the-counter or herbal product
    \item Ask about food and alcohol restrictions for each medicine
    \item Have a Home Medicines Review if you take many medicines
    \item Report new or unusual symptoms after starting a medicine
\end{itemize}

\section*{Sources and Further Reading}
The following reputable sources provide further information for healthcare students and patients seeking reliable, current advice about drug interactions.

\begin{itemize}
    \item NPS MedicineWise. \textit{Drug interactions: what to know.} https://www.nps.org.au/
    \item Healthdirect Australia. \textit{Medicine interactions.} https://www.healthdirect.gov.au/
    \item Therapeutic Goods Administration. \textit{Medicine safety information.} https://www.tga.gov.au/
    \item NHS. \textit{Medicine interactions and safety.} https://www.nhs.uk/
    \item MedlinePlus. \textit{Drug interactions.} https://medlineplus.gov/
    \item Pharmaceutical Society of Australia. \textit{Medication safety resources.} https://www.psa.org.au/
\end{itemize}
